% SHARD-ID: LB-08-LOCAL-RELAXATION
% SHARD-TITLE: Local relaxation, escape obstructions, and kink--magnon wave operators
% SHARD-SOURCES: theory/ace-ld.md; theory/lr-d16.md; theory/ansatz-scattering.md;
%   definitions.md D28 (and, recalled only, D13, D17, D18, D26, D27);
%   claims/CLAIMS.md rows LD-ID, ACE-LD-eps, ACE-LD-sharp, ACE-LD-abs,
%   LR-D16-EDW, LR-D16-NR, ACE-LD-obst-prime, M-ESC-NR, AC-EX, AD3-ex;
%   verdicts theory/verdicts/{ace-ld-r2,ace-ld-r3,ace-ld-r4,lr-d16-r1,lr-d16-r2,
%   blitz-m-esc-nr-r1,memory-index-r3}.md
% SHARD-CLAIMS: LD-ID, ACE-LD-eps, ACE-LD-sharp, ACE-LD-abs, LR-D16-EDW,
%   LR-D16-NR, ACE-LD-obst-prime, M-ESC-NR, AC-EX, AD3-ex; D28

% --- local semantic macros for this shard only (see main.tex for the rest) ---
\newcommand{\Qrel}{\mathcal{Q}}              % relative-charge limit vector
\newcommand{\Nesc}{\mathcal{N}}              % escaped wrong-phase content
\newcommand{\NDW}{N_{\mathrm{DW}}}           % domain-wall number
\providecommand{\lamt}{\widetilde{\lambda}}  % tail decay rate (shared with shard 02)
\newcommand{\dev}{D}                         % on-site tail deviation operator

\section{Local relaxation, escape obstructions, and kink--magnon wave operators}
\label{sec:local-relaxation}

The memory limit law of \cref{sec:memory-index} is not unconditional. It
delivers a quantized ordered wall displacement only on states for which the
windowed charge history relaxes: the three clauses of local relaxation
--- convergence of the fixed-window Ces\`aro states and two-point-measurement
laws, first-moment nondemolition, and first-moment tightness of the family of
window laws --- which we abbreviate (LR1), (LR2), (LR3), and collectively (LR).
This section is the record of what has actually been proved about that
hypothesis and about its enemies.

The material falls into three groups, and it is worth saying in advance which
way each one cuts.

\emph{First}, an exact algebraic identity: the regularised window charge is
nothing but the windowed restriction of the conserved first-moment wall
coordinate. This is unconditional, it costs nothing, and it is the sharpest
single fact in this section, because it converts two apparently
representation-theoretic questions --- can a state have exponentially confined
tails while a charged leg escapes? can such a state carry nonzero ordered
memory? --- into arithmetic. Both answers are negative.

\emph{Second}, a package of theorems around a named exponential-tail
hypothesis (K-TAIL). These say what the confined-charge corner looks like: on
it one has a quantitative window-charge estimate with an explicit
$\lamt^{d_W}$ rate, the tail vacua are forced to have sharp on-site charge,
and --- the sting --- the ordered memory is \emph{exactly zero}. The
hypothesis class is inhabited (the exactly solvable kink model of
\cref{sec:kink-model} sits in it, with numerically measured decay rate), but
it is disjoint from the escaping-leg class in which a nonzero memory could
live.

\emph{Third}, the complementary corner. An unconditional operator inequality
bounds the \emph{number} of phase boundaries in the easy-axis model uniformly
in time, but says nothing about their \emph{length}; the gap between the two
is exactly the residual hypothesis (NR) from which clause (LR3) does follow.
Against this, mean tail transport (M-ESC) is shown to destroy (LR3) outright.
No state realising (M-ESC) is exhibited anywhere in this campaign, so that
result constrains a hypothesis class rather than refuting anything; but
together with the tightness transfer it yields the clean propositional
statement that (M-ESC) and (NR) cannot both hold.

The section closes with the conditional construction of kink--magnon wave
operators from exact ansatz band data, which is the machinery that would
instantiate the channel structure the escape analysis presumes, and with the
local-decay hypothesis it leaves open.

\subsection{Standing notation}
\label{subsec:lr-notation}

The following objects are defined in full in \cref{sec:memory-quantization}
and \cref{sec:memory-index}; they are recalled here so that this section can
be read without turning back.

\begin{definition}[Window observables and the relative window charge]
\label{def:window-charge}
Fix a chain with two $S^z$-vacuum densities $s_\alpha=+s$ at $-\infty$ and
$s_\beta=-s$ at $+\infty$, $s>0$. For a finite window $W=[a,b]\subset\Z$ the
\emph{windowed wall-position observable} is
\begin{equation}
  \wallX_W \;:=\; a-1+\frac{1}{2s}\sum_{x=a}^{b}\bigl(S^z_x+s\bigr),
  \label{eq:wall-observable}
\end{equation}
normalised so that $\wallX_W=m$ on a sharp wall at the bond $(m\mid m+1)$
carrying no other content in $W$. For a cut $c_0\in W$ the \emph{relative
window charge} is $\Qhat_{W,c_0}:=2s(\wallX_W-c_0)$, a local self-adjoint
element of $\alg_W$ with finite spectrum, and $E_W(\cdot)$ denotes its
spectral measure. The \emph{escaped wrong-phase content} outside $W$ is
\begin{equation}
  \Nesc_W \;:=\; \sum_{x<a}\Bigl(\tfrac12-S^z_x\Bigr)
              +\sum_{x>b}\Bigl(S^z_x+\tfrac12\Bigr),
  \label{eq:escaped-content}
\end{equation}
a sum of non-negative diagonal terms, so that $\Nesc_{W'}\le\Nesc_W$ whenever
$W\subseteq W'$.
\end{definition}

\begin{definition}[Tail deviations and window padding]
\label{def:tail-deviations}
For a finite core interval $K=[\ell,r]\subset\Z$ and a site $x\notin K$ put
$\gamma(x):=\alpha$ if $x<\ell$ and $\gamma(x):=\beta$ if $x>r$, and define the
\emph{on-site tail deviation}
\begin{equation}
  \dev_x \;:=\; S^z_x-s_{\gamma(x)},\qquad s_\alpha=+s,\quad s_\beta=-s .
  \label{eq:tail-deviation}
\end{equation}
For a window $W=[a,b]\supseteq K$ the \emph{core-to-edge padding} is
$d_W:=\dist(K,\Z\setminus W)=\min(\ell-a,\;b-r)+1\ge1$. Where a cut $c_0$ is
fixed rather than a core, we write instead the \emph{step density}
$s_{\gamma_{c_0}}(x):=+s$ for $x\le c_0$ and $-s$ for $x>c_0$.
\end{definition}

Throughout, H-MQG(1)--(3) are the standing structural hypotheses of the
general memory-quantization setting of \cref{sec:memory-quantization}: a
compact group $G$ with a covariant family of injective canonical matrix
product state tensors of one common bond dimension and a selected kink sector
$\Kk_{\alpha\beta}$; a common unbroken abelian subgroup with a fixed primitive
on-site charge $S^z$ calibrated by $\omega_\alpha(S^z)=+s=-\omega_\beta(S^z)$,
$s>0$; and a translation-invariant, finite-range, $G$-invariant Hamiltonian
with stationary vacua. D26(INT) is the integrality condition on the on-site
charge, and D17 is the $\ell^1$ kink class.

\subsection{The window charge is the conserved wall coordinate}
\label{subsec:window-charge-identity}

The starting point costs nothing and pays for a great deal. The regularised
window charge, which was introduced as a bookkeeping device for windowed
two-point measurements, is literally the sum of the on-site tail deviations
over the window --- with no offset and no normalisation factor.

\begin{theorem}[The window charge is the windowed sum of tail deviations]
\label{thm:window-charge-identity}
\statusProved\;
Assume H-MQG(1)--(2). For every finite window $W=[a,b]$ with $c_0\in W$,
\begin{equation}
  \Qhat_{W,c_0}\;=\;\sum_{x\in W}\bigl(S^z_x-s_{\gamma_{c_0}}(x)\bigr)
  \qquad\text{exactly.}
  \label{eq:window-charge-identity}
\end{equation}
Consequently, on any state on which the increasing-window limit exists --- in
particular on every state satisfying the exponential-tail hypothesis (K-TAIL)
of \cref{def:k-tail}, where the limit vector
$\Qrel_{c_0}\varphi:=\lim_n\Qhat_{W_n,c_0}\varphi$ exists and is independent
of the exhaustion ---
\begin{equation}
  \Qrel_{c_0}\;=\;\sum_{x\in\Z}\bigl(S^z_x-s_{\gamma_{c_0}}(x)\bigr)
  \;=\;2s\,(X_1-c_0),
  \label{eq:relative-charge-is-X1}
\end{equation}
where $X_1$ is the first-moment wall coordinate of the dynamical dress
(\cref{sec:memory-quantization}), the regularised total magnetisation, which
is exactly conserved. The relative window charge is therefore not a new
object but the vector-valued lift of $X_1$.
\end{theorem}

\begin{scope}
\label{scope:window-charge-identity}
The identity \eqref{eq:window-charge-identity} is unconditional under
H-MQG(1)--(2). The identification \eqref{eq:relative-charge-is-X1} requires in
addition that the increasing-window limit exist on the state in question. What
is constructed is only the limit \emph{vector} $\Qrel_{c_0}\varphi$ on a
(K-TAIL) state, never a sector-wide charge operator: the refutation of the
strong local-additivity claim (\cref{sec:negative-results}) is respected and
not contradicted, its refuting state having polynomially decaying deviations
and therefore lying outside the (K-TAIL) class.
\end{scope}

\begin{proof}[Idea of proof]
The frozen expansion of the windowed observable gives
$\Qhat_{W,c_0}=\sum_{x\in W}S^z_x+s(a+b-1-2c_0)$, and the step density sums to
exactly the negative of that offset,
$\sum_{x\in W}s_{\gamma_{c_0}}(x)=s\,[(c_0-a+1)-(b-c_0)]=-s(a+b-1-2c_0)$. The
two scalars cancel identically; there is nothing left to estimate. Summing the
identity over an exhaustion and comparing with the definition of $X_1$ gives
\eqref{eq:relative-charge-is-X1}.
\end{proof}

\provenance{LD-ID}{\texttt{theory/ace-ld.md} step (1.4).(2.7), THEOREM LD-ID, sub-step (3.1)--(3.2)}{No gate is required: the identity is exact algebra. Independently re-verified in \texttt{theory/verdicts/ace-ld-r3.md} \S0(i) on all $2^{12}$ basis states of the exactly solvable battery over six windows (maximal discrepancy $0.000\mathrm{e}{+00}$), symbolically and on a $201$-site lattice in \texttt{theory/verdicts/ace-ld-r2.md} \S0(i), and reached independently by the parallel lane}{\texttt{ace-ld-r2.md}, \texttt{ace-ld-r3.md}, \texttt{ace-ld-r4.md}; parallel-lane confirmation \texttt{lr-d16-r1.md} \S2 item 1}

The identity has two corollaries, each on its own named hypotheses. The first
closes off an entire class of loopholes.

\begin{corollary}[No escaping charged leg has time-uniform confined tails]
\label{cor:no-escaping-leg}
\statusProvedConditional\;
Assume, in addition to the hypotheses of \cref{thm:window-charge-identity},
that $\Psi_t$ carries an outgoing free leg of charge $q_{\mathrm{leg}}\neq0$
which has left the window $W$ in the sense of the asymptotic-decay clause
D18(AD3), and that the state satisfies (K-TAIL) with \emph{time-independent}
constants $(K,C_K,\lamt)$. Then
\begin{equation}
  \norm{\Qrel_{c_0}\Psi_t-\Qhat_{W,c_0}\Psi_t}
  \;\ge\;\Bigl|\sum_{x\notin W}\omega_{\Psi_t}(\dev_x)\Bigr|
  \;\longrightarrow\;\abs{q_{\mathrm{leg}}}
  \quad(t\to\infty),
  \label{eq:leg-lower-bound}
\end{equation}
while (K-TAIL) bounds the same quantity by $2C_K\lamt^{d_W}/(1-\lamt)$
uniformly in $t$. Choosing $d_W$ so large that
$2C_K\lamt^{d_W}/(1-\lamt)<\abs{q_{\mathrm{leg}}}$ is a contradiction. Hence
no state in the escaping-leg class of D18(AD3) satisfies (K-TAIL) with
time-independent constants, and distinct channel charges $q_L\neq q_T$ are
unavailable within a single conserved-charge sector of D18.
\end{corollary}

\begin{corollary}[On the confined class the ordered memory vanishes]
\label{cor:ordered-memory-zero}
\statusProvedConditional\;
Assume in addition (K-Q) of \cref{def:k-q} at $\varepsilon_Q=0$ together with
time-uniform (K-TAIL) data, and adopt D27's definition of the ordered wall
expectation $\delta x$. Then $q_\varphi$ is a time-independent eigenvalue of
the conserved $\Qrel_{c_0}$, and
\begin{equation}
  \abs{\delta x}\;\le\;s^{-1}\lim_m\frac{2C_K}{1-\lamt}\lamt^{d_{W_m}}\;=\;0 .
  \label{eq:delta-x-zero}
\end{equation}
On that class the ordered memory is exactly zero --- the conservation trap of
the dynamical dress, now in windowed dress.
\end{corollary}

\begin{scope}
\label{scope:ld-id-corollaries}
Neither corollary is claimed outside its stated hypotheses, and nothing is
claimed about which regime any given model occupies.
\Cref{cor:no-escaping-leg} additionally requires D18(AD3) and time-independent
(K-TAIL) constants; \cref{cor:ordered-memory-zero} additionally requires (K-Q)
at $\varepsilon_Q=0$, time-uniform data, and D27's own definition of
$\delta x$.
\end{scope}

\begin{proof}[Idea of proof]
For \cref{cor:no-escaping-leg}: by
\cref{thm:window-charge-identity} the difference
$\Qrel_{c_0}\Psi_t-\Qhat_{W,c_0}\Psi_t$ is exactly the sum of the tail
deviations outside $W$, and Cauchy--Schwarz on unit $\Psi_t$ bounds its norm
below by the modulus of its expectation, which D18(AD3) drives to
$q_{\mathrm{leg}}$. The quantitative witness --- one magnon over a polarised
vacuum --- has escaped tail sum saturating at exactly $2s=1$ at every padding,
with the minimal admissible $C_K$ growing like $\lamt^{-vt}$ (numerically
$10^{+79.8}$ at $t=64$). For \cref{cor:ordered-memory-zero}: the exponential
estimate \eqref{eq:eps-bound} applied on both time wings pins the window-charge
expectation to $q_\varphi$ at every time, and $\delta x$ is by definition
$(2s)^{-1}$ times the limit of the difference of those expectations.
\end{proof}

\provenance{LD-ID}{\texttt{theory/ace-ld.md} step (1.4).(2.7), sub-steps (3.3) and (3.4)}{Corollary (i)'s quantitative witness re-measured in \texttt{theory/verdicts/ace-ld-r3.md} M1(b): $\norm{(\Qhat_{W'}-\Qhat_{W})\varphi}=2s$ at every padding for an escaped leg inside the annulus}{\texttt{ace-ld-r3.md} (F1, M1, m1--m3), \texttt{ace-ld-r4.md}; mechanism reached independently in \texttt{lr-d16-r1.md} M2(b)2}

\subsection{Exponentially confined tails: the \texorpdfstring{$\varepsilon$}{epsilon}-theorem}
\label{subsec:eps-theorem}

We now name the confinement hypothesis precisely and state what it buys.

\begin{definition}[(K-TAIL): confined-core exponential tail decay]
\label{def:k-tail}
A unit vector $\varphi$ satisfies \textbf{(K-TAIL)} if there are a finite
interval $K=[\ell,r]$ and constants $C_K<\infty$, $\lamt\in(0,1)$ such that for
all $x,x'\notin K$ lying on the \emph{same} side of $K$,
\begin{equation}
  \bigl|\omega_\varphi(\dev_x \dev_{x'})\bigr|
  \;\le\; C_K^2\,\lamt^{\dist(x,K)}\,\lamt^{\dist(x',K)} .
  \label{eq:k-tail}
\end{equation}
Taking $x=x'$ gives the on-site bound
$\norm{\dev_x\varphi}^2\le C_K^2\lamt^{2\dist(x,K)}$. Conversely, since each
$\dev_x$ is self-adjoint, Cauchy--Schwarz gives
$\abs{\omega_\varphi(\dev_x\dev_{x'})}
=\abs{\braket{\dev_x\varphi,\dev_{x'}\varphi}}
\le\norm{\dev_x\varphi}\,\norm{\dev_{x'}\varphi}$, so the on-site bound implies
the full same-side clause. \textbf{(K-TAIL) is therefore an on-site tail-decay
hypothesis, not a clustering hypothesis}, and a site-by-site numerical profile
certifies it in full.
\end{definition}

\begin{definition}[(K-Q): sharp relative charge]
\label{def:k-q}
A unit vector $\varphi$ satisfies \textbf{(K-Q)} if there are $q_\varphi\in\R$
and $\varepsilon_Q\ge0$ with
$\norm{\Qrel_{c_0}\varphi-q_\varphi\varphi}\le\varepsilon_Q$, where
$\Qrel_{c_0}\varphi=\lim_n\Qhat_{W_n,c_0}\varphi$ is the relative-charge limit
vector. That limit exists on every (K-TAIL) state and is exhaustion
independent --- this is a lemma, not part of the hypothesis --- and by
\cref{thm:window-charge-identity} it equals $2s(X_1-c_0)$. So (K-Q) is a
sharpness hypothesis on the \emph{conserved} charge, and (K-TAIL) says, up to
the uniform bound below, that no charge lies outside the core.
\end{definition}

\begin{theorem}[Exponential window-charge estimate on confined-tail states]
\label{thm:eps-theorem}
\statusProvedConditional\;
Proved only as the following conditional implication. Assume H-MQG(1)--(2); a
finite core $K=[\ell,r]$; a finite window $W=[a,b]\supseteq K$ with cut
$c_0\in W$ and padding $d_W\ge1$ as in \cref{def:tail-deviations}; and a unit
vector $\varphi$ satisfying (K-TAIL) with constants $C_K,\lamt$ and (K-Q) with
data $q_\varphi,\varepsilon_Q$. Then
\begin{equation}
  \norm{(\Qhat_{W,c_0}-q_\varphi)\varphi}
  \;\le\;\varepsilon_Q+\frac{2C_K}{1-\lamt}\,\lamt^{\,d_W},
  \label{eq:eps-bound}
\end{equation}
and consequently $\dist\bigl(q_\varphi,\;\kappa_{W,c_0}+\Z\bigr)$ is bounded by
the same quantity, where $\kappa_{W,c_0}+\Z$ is the coset containing the
spectrum of the window charge (the finite memory-index theorem of
\cref{sec:memory-index}).
\end{theorem}

Two consequences of (K-TAIL) are forced, and each is itself proved.

\begin{corollary}[Forced consequences of confined tails]
\label{cor:eps-forced}
\statusProvedConditional\;
Under (K-TAIL):
\begin{enumerate}
\item[(i)] In the kink class of D9(a), (K-TAIL) is equivalent to the statement
that no charge lies outside the core, quantitatively
\begin{equation}
  \norm{\Qrel_{c_0}\varphi-\Qhat_{W,c_0}\varphi}
  \;\le\;\frac{2C_K}{1-\lamt}\lamt^{\,d_W}
  \qquad\text{for every finite }W\supseteq K\text{ with }c_0\in W,
  \label{eq:tail-uniform}
\end{equation}
and it forces both tail vacua to have \emph{sharp on-site charge}:
$\omega_\gamma\bigl((S^z_x-s_\gamma)^2\bigr)=0$ at every site, equivalently
$\omega_\gamma\bigl(\Pi_x^{(s_\gamma)}\bigr)=1$ for the on-site spectral
projection of $S^z_x$. Hence $\pm s\in\spec S^z$; in particular a spin-$1$
chain with calibration $s=1/2$ admits no (K-TAIL) kink state. See
\cref{thm:sharp-charge}.
\item[(ii)] With time-independent data $(K,C_K,\lamt)$, (K-TAIL) is
\emph{jointly unsatisfiable} with the clause of D18(AD3) that the free leg
charge leaves the window. No state carrying an escaping leg of nonzero charge
satisfies it with time-independent constants; $q_L\neq q_T$ is unavailable
within one conserved-charge sector of D18; and on states that do satisfy it
with $\varepsilon_Q=0$ and time-uniform data the ordered wall displacement is
$\delta x=0$.
\end{enumerate}
\end{corollary}

\begin{scope}
\label{scope:eps-theorem}
The following are \emph{not} claimed. That the $\varepsilon$-weakened
local-decay statement supplies the cross-term step of the wave-operator
construction of \cref{subsec:wave-operators} --- it does not, because the
hypothesis fails on that entire class; the earlier claim to the contrary is
retracted. The exact fixed-window local-decay hypothesis (AD3-ex) of
\cref{conj:ad3-ex}. The weak-$*$ clause of D18(AD3) or the reduction clause of
the spectral memory-index statement. The $\chi=1$ product form of the tail
vacua under H-MQG alone (see \cref{ref:chi-one}). Any statement about generic
injective matrix product state vacua --- sharp on-site charge excludes every
vacuum with nonzero on-site charge variance, and that is the scope boundary of
this theorem. Any infinite-volume proof of (K-TAIL) or (K-Q) on any model.
The existence of a sector-wide conserved charge operator: only the limit
vector on a (K-TAIL) state is constructed, the strong local-additivity claim
is refuted and is not contradicted, and the infinite-volume route to the
time-uniformity of (K-Q) remains a conjecture.
\end{scope}

\begin{proof}[Idea of proof]
Two lemmas. The first is a window-difference identity: for
$W=[a,b]\subseteq W'=[a',b']$ both containing the cut,
\begin{equation}
  \Qhat_{W',c_0}-\Qhat_{W,c_0}
  \;=\;\sum_{x=a'}^{a-1}\bigl(S^z_x-s\bigr)
      +\sum_{x=b+1}^{b'}\bigl(S^z_x+s\bigr),
  \label{eq:window-difference}
\end{equation}
so the difference is exactly the sum of the on-site deviations over the
annulus, the cut dependence cancelling because $c_0$ enters both terms with the
same coefficient. The second is a tail-sum estimate: for a finite one-sided set
$F$ at distance $d$ from $K$, expanding the square and using (K-TAIL) on every
pair gives
$\norm{\sum_{x\in F}\dev_x\varphi}\le C_K\lamt^{d}/(1-\lamt)$, because the
distinct sites of a one-sided $F$ have pairwise distinct integer distances to
$K$, all at least $d$, so the geometric series sums. Applying this to the two
annuli in \eqref{eq:window-difference} shows that
$\{\Qhat_{W_n,c_0}\varphi\}$ is Cauchy along any exhaustion, that interleaving
two exhaustions identifies the limits, and that letting $W'\uparrow\Z$ gives
\eqref{eq:tail-uniform}. The theorem is then the triangle inequality on
$(\Qhat_{W,c_0}-q_\varphi)\varphi
=(\Qhat_{W,c_0}\varphi-\Qrel_{c_0}\varphi)+(\Qrel_{c_0}\varphi-q_\varphi\varphi)$,
and the coset location follows from
$\norm{(\Qhat_{W,c_0}-q_\varphi)\varphi}^2=\int\abs{\lambda-q_\varphi}^2\,
d\mu_\varphi(\lambda)\ge\dist(q_\varphi,\spec\Qhat_{W,c_0})^2$ together with
spectral permanence.
\end{proof}

\begin{remark}[Why the time-uniform hypothesis is honest]
\label{rem:kq-propagation}
(K-Q) constrains the \emph{conserved} charge, and that is what makes the
time-uniform version of \cref{thm:eps-theorem} an assumption about the initial
data rather than a disguised assumption of the conclusion. If the realization
carries a self-adjoint $\widehat{\Qrel}$ commuting with the dynamics and
agreeing with the limit vector on (K-TAIL) states, then
$\norm{(\widehat{\Qrel}-q)\varphi_t}=\norm{(\widehat{\Qrel}-q)\varphi_0}$, so
(K-Q) at $t=0$ propagates with the same $q$ and $\varepsilon_Q$. The window
charge $\Qhat_{W,c_0}$ is \emph{not} conserved --- that is the whole content of
the conservation trap --- and no such propagation is claimed for it: the
estimate converts conserved-charge sharpness plus instantaneous tail decay into
a window statement separately at each time.
\end{remark}

\begin{remark}[The hypothesis class is inhabited]
\label{rem:d16-instantiation}
The exactly solvable easy-axis kink model of \cref{sec:kink-model} sits inside
the class. Its exact zero-energy sector states satisfy (K-Q) exactly, and they
satisfy (K-TAIL) at rate
\begin{equation}
  \lamt \;=\; q \;=\; \Delta-\sqrt{\Delta^2-1},
  \label{eq:d16-rate}
\end{equation}
independently verified by full $2^L$ exact diagonalisation at $L=12$: the
measured on-site defects are
$\norm{\dev_x\varphi}=4.356\times10^{-2},\;9.092\times10^{-3},\;
1.898\times10^{-3},\;3.960\times10^{-4},\;8.266\times10^{-5}$, with per-site
ratios $0.20871$ on \emph{both} sides against $q=0.208712$, and $C_K=0.2087$,
in the $\Z_2$ image of the frozen orientation. The checker's conclusion gates
independently certify the $\lamt^{d_W}$ rate of \eqref{eq:eps-bound}, with
defects from $4.45\times10^{-2}$ down to $4.05\times10^{-4}$ and ratios
$0.2089$, $0.2087$, $0.2083$ against the same $q$.
\end{remark}

\provenance{ACE-LD-eps}{\texttt{theory/ace-ld.md} step (1.4) (THEOREM ACE-LD-$\varepsilon$), with lemmas at (1.4).(2.1)--(2.3) and the forced consequences at (1.4).(2.7), (2.9)}{\texttt{theory/checks/ace\_ld\_check.py} gate LD-C7 (green exit $0$; \texttt{--red} exit $1$, RED-OK $23/23$). Gates (b)--(c) certify the conclusion's $\lamt^{d_W}$ rate; gate (e) certifies the hypothesis (K-TAIL) itself site by site, and by Cauchy--Schwarz on the self-adjoint $\dev_x$ the on-site clause is equivalent to the full same-side clause, so (e) certifies (K-TAIL) in full. No gate bears on (K-Q), on the infinite-volume statement, or on the sharp-charge corollary, which is exact algebra}{\texttt{ace-ld-r3.md} \S8(A) (adjudicated scoping), \texttt{ace-ld-r2.md}, \texttt{ace-ld-r4.md} \S0(iv)}

\subsection{Sharp on-site charge, and the refuted product-vacuum reading}
\label{subsec:sharp-charge}

The forced consequence (i) above deserves its own statement, because it is the
cleanest description of exactly which corner of the theory the
$\varepsilon$-theorem occupies --- and because the natural strengthening of it
is false.

\begin{theorem}[Confined tails force zero on-site charge variance]
\label{thm:sharp-charge}
\statusProved\;
Assume H-MQG(1)--(2), and let $\varphi$ be a unit vector whose state
$\omega_\varphi$ lies in the kink sector $\Kk_{\alpha\beta}$ of D9(a) and
satisfies (K-TAIL) with core $K$ and constants $C_K,\lamt$. Then both tail
vacua have \emph{sharp on-site charge}: for $\gamma\in\{\alpha,\beta\}$ and
every site $x$,
\begin{equation}
  \omega_\gamma\bigl((S^z_x-s_\gamma)^2\bigr)=0,
  \qquad\text{equivalently}\qquad
  \omega_\gamma\bigl(\Pi_x^{(s_\gamma)}\bigr)=1,
  \label{eq:sharp-charge}
\end{equation}
where $\Pi_x^{(m)}$ is the on-site spectral projection of $S^z_x$ at
eigenvalue $m$. Hence the tail density is an on-site \emph{eigenvalue}:
$s_\alpha=+s\in\spec S^z$ and $s_\beta=-s\in\spec S^z$.
\end{theorem}

\begin{corollary}[An arithmetic exclusion]
\label{cor:spin-one-exclusion}
\statusProved\;
A spin-$1$ chain with calibration $s=1/2$ admits \emph{no} (K-TAIL) kink
state, since $\pm1/2\notin\spec S^z=\{-1,0,1\}$.
\end{corollary}

\begin{scope}
\label{scope:sharp-charge}
Sharp on-site charge means zero on-site charge variance and therefore excludes
every vacuum with nonzero on-site charge variance, in particular generic
injective matrix product state vacua. It is \emph{not} claimed that the tail
vacua are fully polarised $S^z$-product states; see \cref{ref:chi-one}. That
stronger form needs $s_\gamma$ to be a \emph{simple} eigenvalue of $S^z$ ---
true in the standard spin-$S$ register with on-site dimension $d=2s+1$, and
true in the exactly solvable model, but not implied by H-MQG(1)--(2), by
covariance, or by D26(INT), all of which admit degenerate on-site charge
eigenspaces. The gloss $d=2s+1$ in the notation table is explicitly fenced as
the fully polarised special case and is not a constraint.
\end{scope}

\begin{proof}[Idea of proof]
Every input is named. D9(a) itself supplies the weak-$*$ tail relaxation
$\omega_\varphi(\tau_x(O))\to\omega_\gamma(O)$ for local $O$, so no extra
relaxation hypothesis is needed. (K-TAIL) at $x=x'$ gives
$\omega_\varphi(\dev_x^2)\le C_K^2\lamt^{2\dist(x,K)}\to0$; translating along
the tail and passing to the limit gives
$\omega_\gamma\bigl((S^z_x-s_\gamma)^2\bigr)=0$. Cauchy--Schwarz on the
self-adjoint $S^z_x-s_\gamma$ then pins the on-site spectral weights, and
translation invariance of $\omega_\gamma$ makes the conclusion site
independent.
\end{proof}

\begin{refutedclaim}[The tail vacua are fully polarised product states]
\label{ref:chi-one}
\statusRefuted\;
\emph{Withdrawn statement.} That \cref{thm:sharp-charge} concludes further
$\omega_\gamma=\bigotimes_x\ket{s_\gamma}\bra{s_\gamma}$, a bond-dimension-one
$S^z$-product state.

\emph{Why it fails.} Sharp on-site charge names a unique on-site \emph{vector}
only if $\dim\ker(S^z-s_\gamma)=1$, and no hypothesis in force supplies that
simplicity. Explicit counterexample: take the on-site space $\C^2\otimes\C^3$
with $S^z:=\sigma^z/2\otimes\mathbb{1}_3$, so that
$\spec S^z=\{\pm1/2\}$ with each eigenvalue of multiplicity $3$ and
$e^{2\pi iS^z}=-\mathbb{1}$, whence D26(INT) holds with $\kappa=1/2$. Put
\begin{equation}
  \omega_\alpha=\bigotimes_x\bigl(\ket{\uparrow}\bra{\uparrow}\otimes
  \varrho_{\mathrm{AKLT}}\bigr),\qquad
  \omega_\beta=\bigotimes_x\bigl(\ket{\downarrow}\bra{\downarrow}\otimes
  \varrho_{\mathrm{AKLT}}\bigr).
  \label{eq:chi-two-counterexample}
\end{equation}
These are injective matrix product states of bond dimension $\chi=2$ (bond
Schmidt spectrum $(1/2,1/2)$, transfer gap $1/3$), covariant for
$G=U(1)\rtimes\Z_2$, with $\omega_\alpha(S^z)=+1/2=s=-\omega_\beta(S^z)$. The
sharp domain wall in the $\C^2$ factor lies in $\Kk_{\alpha\beta}$ and has
$\dev_x\varphi=0$ at \emph{every} site, so it satisfies (K-TAIL) with
$C_K=0$ --- while its tail vacua are not product states.

\emph{Surviving weaker statement.} \Cref{thm:sharp-charge}: zero on-site charge
variance, hence $\pm s\in\spec S^z$. When $s_\gamma$ \emph{is} a simple
eigenvalue of $S^z$ --- the standard spin-$S$ register, and the exactly
solvable model of \cref{sec:kink-model} --- the sharp-charge conclusion does
give $\omega_\gamma=\bigotimes_x\ket{s_\gamma}\bra{s_\gamma}$ and $\chi=1$;
that conditional clause survives as a corollary, not as the theorem.
\end{refutedclaim}

\provenance{ACE-LD-sharp}{\texttt{theory/ace-ld.md} step (1.4).(2.9)(b)}{\texttt{theory/checks/ace\_ld\_check.py} gate LD-C7(e) certifies the \emph{hypothesis} (K-TAIL) on the exactly solvable instance (per-site ratios $0.20871$ on both sides, $C_K=0.209$, red-armed by \texttt{--red-c7-orientation} at $C_K=2525$); the conclusion is exact algebra and needs no gate. The counterexample to the product form is verified numerically}{\texttt{ace-ld-r3.md} F1 and F1(b) (the refutation, adopted), \texttt{ace-ld-r4.md} (restatement)}

\subsection{The abstract mechanism: spectral diagonality from first-moment escape}
\label{subsec:spectral-diagonality}

Underneath the model-specific analysis sits a short, entirely abstract lemma.
It is worth isolating because it shows exactly what local decay of the channel
charges buys: nothing about the model, but a clean identification of channel
projections with spectral projections of the window charge, with a constant
uniform in time.

\begin{definition}[The abstract escape package (A1)--(A5)]
\label{def:a1-a5}
\begin{enumerate}
\item[(A1)] A Hilbert space $\Hilb$, a strongly continuous unitary group
$e^{-itH}$, a unit vector $\Psi$, and $\Psi_t:=e^{-itH}\Psi$.
\item[(A2)] A bounded self-adjoint operator $\Qhat$ on $\Hilb$, time
independent (Schr\"odinger picture), whose spectrum is a \emph{finite} set
contained in a single coset $\kappa+\Z$, $\kappa\in[0,1)$. In the instance
$\Qhat=\Qhat_{W,c_0}$ at a fixed finite window, and the coset containment is
the finite memory-index theorem of \cref{sec:memory-index}, cited and not
re-proved.
\item[(A3)] A finite family $\{P_{\mathrm{ch}}\}_{\mathrm{ch}=1}^{n}$ of
mutually orthogonal projections, each commuting with $e^{-itH}$, with
$\sum_{\mathrm{ch}}P_{\mathrm{ch}}\Psi=\Psi$.
\item[(A4)] Real \emph{channel charges}
$\{q_{\mathrm{ch}}\}_{\mathrm{ch}=1}^{n}$ and first-moment escape:
$\varepsilon_{\mathrm{ch}}(t):=
\norm{(\Qhat-q_{\mathrm{ch}})P_{\mathrm{ch}}\Psi_t}\to0$ as $t\to+\infty$ for
every channel.
\item[(A5)] The $q_{\mathrm{ch}}$ are pairwise distinct.
\end{enumerate}
\end{definition}

\begin{lemma}[Channel projections become spectral projections]
\label{thm:spectral-diagonality}
\statusProved\;
Assume (A1)--(A5). With $E(\cdot)$ the spectral measure of $\Qhat$ and
$d_{\mathrm{ch}}:=\dist\bigl(q_{\mathrm{ch}},\spec\Qhat\setminus
\{q_{\mathrm{ch}}\}\bigr)>0$,
\begin{equation}
  \norm{E(\{q_{\mathrm{ch}}\})\Psi_t-P_{\mathrm{ch}}\Psi_t}
  \;\le\;\sum_{\mathrm{ch}'=1}^{n}d_{\mathrm{ch}'}^{-1}
  \varepsilon_{\mathrm{ch}'}(t)\;\longrightarrow\;0
  \qquad\text{for every channel,}
  \label{eq:spectral-diagonality}
\end{equation}
the constant $\sum_{\mathrm{ch}'}d_{\mathrm{ch}'}^{-1}$ being independent of
$t$. If moreover $q_{\mathrm{ch}}\in\kappa+\Z$ then $d_{\mathrm{ch}}\ge1$, so
under (A2) the constant is at most $n$ for coset-valued charges. The lemma
holds for $N$ channels and a general on-site spin.
\end{lemma}

\begin{lemma}[Distinct channel charges are necessary]
\label{thm:distinct-charges-necessary}
\statusProved\;
Assume (A1)--(A4) with $n=2$ and $q_1=q_2=:q$, and
$w_i:=\liminf_{t\to+\infty}\norm{P_i\Psi_t}>0$ for \emph{both} $i=1,2$. Then
for \emph{every} $q'\in\R$,
$\liminf_{t\to+\infty}\norm{E(\{q'\})\Psi_t-P_1\Psi_t}>0$; the diagonality
display fails for channel $1$ for every charge assignment, and by symmetry for
channel $2$. Hypothesis (A5) is therefore not an artefact of the proof.
\end{lemma}

\begin{scope}
\label{scope:spectral-diagonality}
Nothing is claimed about any concrete model, no instantiation of (A4) is
supplied, and no sharpness of the constant beyond the single spectral gap is
asserted.
\end{scope}

\begin{proof}[Idea of proof]
The engine is a one-line gap estimate: if $\Qhat$ has finite spectrum $S$ and
$d_q:=\dist(q,S\setminus\{q\})>0$, then $\Qhat$ restricted to
$\ran(1-E(\{q\}))$ has spectrum $S\setminus\{q\}$, so $(\Qhat-q)$ is invertible
there with inverse norm $d_q^{-1}$, giving
\begin{equation}
  \norm{(1-E(\{q\}))\phi}\;\le\;d_q^{-1}\norm{(\Qhat-q)\phi}
  \qquad\text{for every }\phi\in\Hilb .
  \label{eq:gap-estimate}
\end{equation}
Because the projections resolve $\Psi$ and commute with the dynamics, they
resolve $\Psi_t$ for all $t$, and
$E(\{q_{\mathrm{ch}}\})\Psi_t-P_{\mathrm{ch}}\Psi_t
=-(1-E(\{q_{\mathrm{ch}}\}))P_{\mathrm{ch}}\Psi_t
+\sum_{\mathrm{ch}'\neq\mathrm{ch}}E(\{q_{\mathrm{ch}}\})
P_{\mathrm{ch}'}\Psi_t$. The diagonal term is bounded by
\eqref{eq:gap-estimate} directly. For a cross term, distinctness of the charges
makes the two spectral projections orthogonal, so
$E(\{q_{\mathrm{ch}}\})P_{\mathrm{ch}'}\Psi_t
=E(\{q_{\mathrm{ch}}\})(1-E(\{q_{\mathrm{ch}'}\}))P_{\mathrm{ch}'}\Psi_t$ and
\eqref{eq:gap-estimate} applies again at $q_{\mathrm{ch}'}$. Time uniformity is
free: $\spec\Qhat$ is a property of the fixed Schr\"odinger-picture operator,
and only $\Psi_t$ moves. Two distinct points of one coset of $\Z$ differ by a
nonzero integer, which is where the coset containment enters.
\end{proof}

\provenance{ACE-LD-abs (with ACE-LD-nec)}{\texttt{theory/ace-ld.md} steps (1.1)--(1.3): the gap estimate at (1.1), the assembly at (1.2), the necessity of distinct charges at (1.3)}{\texttt{theory/checks/ace\_ld\_check.py} gates LD-C1--C4 (green exit $0$; \texttt{--red} exit $1$, $23$ mutation modes, exit paths logged; every LD-C1 sub-gate armed, including orthogonality and resolution)}{\texttt{ace-ld-r2.md} (M4/m4 fixed), \texttt{ace-ld-r3.md}, \texttt{ace-ld-r4.md}}

\subsection{Energy bounds the number of phase boundaries, not their length}
\label{subsec:edw}

We turn to the complementary corner: the easy-axis kink model of
\cref{sec:kink-model}, where one would like to \emph{prove} local relaxation
rather than assume it. There is an unconditional operator inequality, and it
is instructive precisely because of what it fails to give.

\begin{theorem}[Energy dominates the domain-wall number]
\label{thm:edw}
\statusProved\;
In the easy-axis model, in the two-site basis
$(\ket{\uparrow\uparrow},\ket{\uparrow\downarrow},
\ket{\downarrow\uparrow},\ket{\downarrow\downarrow})$, the bond term
$h^{\mathrm{XXZ}}_{x,x+1}$ annihilates $\ket{\uparrow\uparrow}$ and
$\ket{\downarrow\downarrow}$ and equals $(J\Delta/2)I-(J/2)\sigma^x$ on the
domain-wall block, with eigenvalues $(J/2)(\Delta\mp1)$. Hence, with
$P^{\mathrm{DW}}_x$ the projection onto
$\mathrm{span}\{\ket{\uparrow\downarrow},\ket{\downarrow\uparrow}\}_{x,x+1}$
and $\NDW:=\sum_xP^{\mathrm{DW}}_x$ the domain-wall \emph{number},
\begin{equation}
  h^{\mathrm{XXZ}}_{x,x+1}\;\succeq\;\tfrac{J}{2}(\Delta-1)P^{\mathrm{DW}}_x,
  \qquad
  H_{\mathrm{XXZ}}\;\succeq\;\tfrac{J}{2}(\Delta-1)\NDW\;\succeq\;0,
  \label{eq:edw-operator}
\end{equation}
the sum being a quadratic-form inequality on the finite-deviation core.
Consequently, \emph{for any state $\Psi$ whose energy
$E_0:=\braket{\Psi,H_{\mathrm{XXZ}}\Psi}$ is finite and is conserved by the
dynamics},
\begin{equation}
  \braket{\Psi,\alpha_t(\NDW)\Psi}\;\le\;\frac{2E_0}{J(\Delta-1)}
  \qquad\text{for every }t .
  \label{eq:edw-bound}
\end{equation}
\end{theorem}

\begin{scope}
\label{scope:edw}
This bounds the \textbf{number} of phase boundaries uniformly in time and
\textbf{nothing about their length}: in one dimension a wrong-phase block of
any length carries the same $O(1)$ Ising cost. It therefore does \emph{not}
give clause (LR3), and it gives neither clause of the hypothesis (NR) of
\cref{def:nr}. Finiteness and conservation of $E_0$ are stated as hypotheses
on the selected state, not derived: the proof step asserts both without proof,
and the operator-hygiene paragraph of the shard covers only the diagonal
operators, not $H_{\mathrm{XXZ}}$. Both are expected to hold on an $\ell^1$
packet of D17 --- $P^{\mathrm{DW}}_x\preceq n_x+n_{x+1}$ and
$0\preceq h_x\preceq(J/2)(\Delta+1)P^{\mathrm{DW}}_x$ give $E_0<\infty$ from
the $\ell^1$ clause --- but that derivation is not in the shard. Finally, the
bond inequality and its sum are exact in infinite volume \emph{and} on any
finite open chain; a finite-chain certificate must use the
$H_{\mathrm{XXZ}}$ propagator, because $H_{\mathrm{kink}}$ and
$H_{\mathrm{XXZ}}$ generate the same derivation only in infinite volume and
differ at finite $N$ by the boundary term
$\tfrac{J}{2}\sqrt{\Delta^2-1}\,(S^z_1-S^z_N)$.
\end{scope}

\begin{proof}[Idea of proof]
The bond term is block diagonal: it kills the two aligned configurations, and
on the two-dimensional domain-wall block it is
$(J\Delta/2)I-(J/2)\sigma^x$, whose eigenvalues are $(J/2)(\Delta\mp1)$. For
$\Delta>1$ the smaller one is positive, so the bond term dominates
$(J/2)(\Delta-1)$ times the domain-wall projection. Summing bondwise gives
\eqref{eq:edw-operator}, and taking the expectation in a state of conserved
finite energy gives \eqref{eq:edw-bound}. As a sanity check, the bare kink has
$E_0=C_K=\tfrac{J}{2}\sqrt{\Delta^2-1}$, so
$\braket{\NDW}\le\sqrt{(\Delta+1)/(\Delta-1)}$, which tends to $1$ as
$\Delta\to\infty$ (a single sharp wall) and diverges as $\Delta\to1^+$ (the
kink delocalises) --- both as they must.
\end{proof}

\provenance{LR-D16-EDW}{\texttt{theory/lr-d16.md} step (1.5).(2.6), named computation LRD-EDW}{\texttt{theory/checks/lr\_d16\_check.py} gate C4: (a) the bond inequality, minimum eigenvalue $-0.0$; (b) $H_{\mathrm{XXZ}}$ energy drift $1.11\mathrm{e}{-14}$ and $\max\braket{\NDW}=3.053\le4.753$; (c) finite-chain calibration $\braket{K|H_{\mathrm{XXZ}}|K}=1.1456439081$ against $C_K=1.1456439237$. Green exit $0$; the mutations \texttt{--red c4-wrong-gap}, \texttt{--red c4-kink-propagator} (caught by the energy-drift gate at $1.007$) and \texttt{--red c4-sharp-calibration} each exit $1$ on exactly that row}{\texttt{lr-d16-r1.md}, \texttt{lr-d16-r2.md}}

\subsection{Tightness transfer: from bounded escaped content to first-moment tightness}
\label{subsec:nr-transfer}

The gap left open by \cref{thm:edw} --- block length rather than block number
--- is named, once, as a hypothesis, and clause (LR3) is derived from it.

\begin{definition}[(NR): no unbounded escaped content, two clauses]
\label{def:nr}
Let $W_1$ be the smallest window of the padded exhaustion, $\Psi$ the selected
vector, $E_{W,t}$ the spectral resolution of the Heisenberg translate of
$\Qhat_{W,c_0}$, and $\Nesc_{W}$ as in \eqref{eq:escaped-content}. Then (NR)
consists of
\begin{align}
  \sup_{t\in\R}\;\braket{\Psi,\alpha_t(\Nesc_{W_1})^2\Psi}&<\infty,
  \label{eq:nr-clause-1}\\
  \sup_m\;\sup_{t_-<t_+}\;\sum_q
  \norm{\alpha_{t_+}(\Nesc_{W_1})E_{W_m,t_-}(\{q\})\Psi}^2&<\infty,
  \label{eq:nr-clause-2}
\end{align}
and $S_{\mathrm{NR}}$ denotes twice the supremum of the sum of the two
quantities.
\end{definition}

\begin{theorem}[Bounded escaped content implies first-moment tightness]
\label{thm:nr-transfer}
\statusSketch\;
In the standing register (H1)--(H6) of the easy-axis setting, with clause
(LR1), the two-clause hypothesis (NR) implies clause (LR3) for the same padded
exhaustion and the same fixed-window laws, with the explicit tail bound
\begin{equation}
  \sup_m\sum_{\abs{\nu}>M}(1+\abs{\nu})\,p_{W_m}(\nu)
  \;\le\;\frac{2S_{\mathrm{NR}}}{M}\;\longrightarrow\;0 .
  \label{eq:nr-tail-bound}
\end{equation}
\end{theorem}

\begin{scope}
\label{scope:nr-transfer}
The status is that of an exact conditional proof with an independent critic
re-derivation, whose promotion in the claims graph is still pending; the
implication itself is what is claimed, not its hypothesis. Specifically, the
following are \emph{not} claimed: (NR) itself; clause (LR2); the optional
weak-convergence clause of (LR3); and either clause of (NR) as a consequence of
\cref{thm:edw}, which controls wall number, not block length.
\end{scope}

\begin{proof}[What is proved, and how]
Four exact steps and one hypothesis. (1) An exact second-moment identity: with
$\Delta\Qhat_W:=\Qhat_{W,c_0}(t_+)-\Qhat_{W,c_0}(t_-)$ and
$\mathcal{D}_{W,t_-}$ the pinching by the spectral projections at $t_-$,
$\sum_\nu\nu^2p_{W;t_-,t_+}(\nu)
=\braket{\Psi,\mathcal{D}_{W,t_-}\bigl((\Delta\Qhat_W)^2\bigr)\Psi}$.
(2) Chebyshev: for any probability $p$ on $\Z$ and $M\ge1$,
$\sum_{\abs{\nu}>M}(1+\abs{\nu})p(\nu)\le
2\sum_{\abs{\nu}>M}\abs{\nu}p(\nu)\le(2/M)\sum_\nu\nu^2p(\nu)$, so a uniform
second moment is exactly first-moment tightness.
(3) A splitting onto the two window edges: $\Delta\Qhat_W$ equals the
difference of the two Heisenberg translates of the complementary charge (the
scalar cancels), and $2X^2+2Y^2-(X-Y)^2\succeq0$ with positivity of the
pinching gives
\begin{equation}
  \sum_\nu\nu^2p_{W;t_-,t_+}
  \;\le\;2\braket{\Psi,\alpha_{t_-}(\Qhat_{W^c})^2\Psi}
  +2\sum_q\norm{\alpha_{t_+}(\Qhat_{W^c})\varphi_q}^2 .
  \label{eq:second-moment-split}
\end{equation}
(4) A monotonicity lemma: $\Nesc_{W'}\le\Nesc_W$ for $W\subseteq W'$ and
$\abs{\Qhat_{W^c,c_0}}\le\Nesc_W$, so both terms may be majorised with the
\emph{fixed} observable $\Nesc_{W_1}$. That is (NR). The remaining
$m$-dependence sits in the dephased state, not in the observable --- positivity
of a pinching does not make its expectation of a non-commuting positive
observable monotone --- which is exactly why the second clause of (NR) retains
its supremum over $m$. Convexity of the double Ces\`aro average and Fatou's
lemma transport the uniform bound to the fixed-window limit laws.

\emph{What is missing.} (NR) itself. It is genuine and non-circular: it
controls supremum-in-time second moments of the state and of the dephased
states, whereas (LR3) controls only tails of double-Ces\`aro limit laws, and
(LR3) does not recover either clause. The honest reading is that
\cref{thm:edw} bounds the number of wrong-phase blocks unconditionally, (NR)
is the statement that their length also stays bounded, and this theorem is the
bridge from the latter to (LR3).
\end{proof}

\begin{remark}[Why (NR) is the right hypothesis here]
\label{rem:flatness}
In the exactly solvable model the closed span of the zero-energy kink family is
annihilated by the kink Hamiltonian, so the dynamics acts trivially on it up to
a constant phase: every state in that closed span is exactly stationary. This
flatness is what makes an assumption about escaped content, rather than about
transport, the natural residual input. Exhaustion of the kernel by that family
is a separate conjecture and is not used here.
\end{remark}

\provenance{LR-D16-NR}{\texttt{theory/lr-d16.md} step (1.5).(2.1)--(2.7), THEOREM LRD-3; independently re-derived in \texttt{theory/verdicts/lr-d16-r2.md} \S4 item 3}{\texttt{theory/checks/lr\_d16\_check.py} gates C3(a)--C3(c) test finite identities and majorants only; no finite checker bears on the asymptotic implication}{\texttt{blitz-m-esc-nr-r1.md} (0F/0M/0m), \texttt{lr-d16-r1.md}, \texttt{lr-d16-r2.md}}

\subsection{The escape dichotomy: mean transport destroys tightness}
\label{subsec:escape-dichotomy}

The opposite hypothesis --- that the kink's mean position transits the window
at a rate proportional to the window size --- is incompatible with local
relaxation, and the argument needs remarkably little.

\begin{definition}[(M-ESC): mean tail transport]
\label{def:m-esc}
In the setting of D27 (a D17 vector $\Psi$, a cut $c_0$, a padded exhaustion
$W_m=[a_m,b_m]\uparrow\Z$ containing $c_0$) with the fixed-window Ces\`aro
states $\omega^\pm_{W_m}$, the hypothesis \textbf{(M-ESC)} is
\begin{equation}
  \theta_{\mathrm{tr}}\;:=\;\liminf_m
  \frac{\bigl|\omega^+_{W_m}(\wallX_{W_m})
        -\omega^-_{W_m}(\wallX_{W_m})\bigr|}{\abs{W_m}}\;>\;0 .
  \label{eq:m-esc}
\end{equation}
\end{definition}

\begin{theorem}[Mean escape is incompatible with first-moment tightness]
\label{thm:escape-obstruction}
\statusProvedConditional\;
Assume H-MQG(1)--(3), D26(INT), and the D27 setting with clauses (LR1)--(LR2)
but \emph{not} (LR3); assume (M-ESC). Then
\begin{enumerate}
\item[(a)] $\bigl|\sum_\nu\nu\,p_{W_m}(\nu)\bigr|
\ge2s\,\theta_{\mathrm{tr}}\abs{W_m}(1-o(1))\to\infty$;
\item[(b)] clause (LR3) \textbf{fails}:
$\sup_m\sum_{\abs{\nu}>M}(1+\abs{\nu})p_{W_m}(\nu)=\infty$ for every $M$;
\item[(c)] the ordered wall expectation $\delta x$ of D27 is undefined along
that exhaustion, its defining differences diverging linearly at rate at least
$\theta_{\mathrm{tr}}$ per window site.
\end{enumerate}
\textbf{Positive form}, which is not a contrapositive: the same step gives
directly that on any state satisfying (LR) the mean wall transport is
\emph{uniformly bounded},
\begin{equation}
  \sup_m\bigl|\omega^+_{W_m}(\wallX_{W_m})
  -\omega^-_{W_m}(\wallX_{W_m})\bigr|\;\le\;\frac{M_0+1}{2s},
  \label{eq:positive-transport-bound}
\end{equation}
where $M_0$ is any threshold with
$\sup_m\sum_{\abs{\nu}>M_0}(1+\abs{\nu})p_{W_m}(\nu)\le1$; in particular
$\theta_{\mathrm{tr}}=0$.
\end{theorem}

\begin{scope}
\label{scope:escape-obstruction}
\textbf{No model or state realising (M-ESC) is exhibited anywhere in this
corpus} --- independently confirmed by the parallel lane. The theorem
therefore constrains a hypothesis class, and until a witness exists it makes no
clause of any class theorem necessary.

Disclosure: boundedness of the mean transport under (LR3) is the quantitative
form of D27's own existence corollary, already frozen in D27's definition
paragraph. The genuine addition is the $M_0$ tightness route, and it is a
strict strengthening: D27's corollary delivers convergence along the full
sequence only under D27's optional convenience clause, and otherwise only along
an (LR3) subsequence, whereas \eqref{eq:positive-transport-bound} is a supremum
over \emph{all} $m$, with no subsequence, no convenience clause, and no passage
through weak convergence.

Nothing is claimed about the exactly solvable model or any confined-kink class.
There $\theta_{\mathrm{tr}}=0$, but only on named conditionals: every state in
the closed span of the zero-energy kink family is exactly stationary (proved),
while exhaustion of the kernel by that family is a conjecture; and the
kink--magnon bound $\abs{\delta x}\le\braket{N_T}/s$ is the conclusion of the
general memory-quantization theorem, which is proved \emph{conditional} on the
full asymptotic-decay package D18(AD1)--(AD4) for that vector --- open for this
model, since (AD3-ex) is a conjecture (\cref{conj:ad3-ex}). Note that D18(AD3)
is jointly unsatisfiable with (K-TAIL) (\cref{cor:eps-forced}(ii)), so this
$\theta_{\mathrm{tr}}=0$ fence and the $\varepsilon$-theorem of
\cref{subsec:eps-theorem} concern \emph{disjoint} state classes: a
kink--magnon $\ell^1$ packet under the asymptotic-decay package, versus a
magnon-free confined-tail sector state. Also not claimed: any statement about
the packet class of D28 as such, which contains $\theta_{\mathrm{tr}}=0$
members; mass-defectiveness of weak limits (first-moment divergence only); and
the two-atom law as a theorem --- it is numerically exhibited only.
\end{scope}

\begin{proof}[Idea of proof]
Three ingredients, all cited rather than re-proved. First, each $p_{W_m}$ is a
probability on $\Z$: at fixed $m$ the law is supported on finitely many $\nu$,
because $\abs{\nu}$ is bounded by the diameter of the finite set
$\spec\Qhat_{W_m,c_0}$, so no mass escapes in the Ces\`aro limit. Second, the
exact first-moment identity, valid at every fixed $m$,
\begin{equation}
  \sum_\nu\nu\,p_{W_m}(\nu)
  \;=\;-2s\bigl[\omega^+_{W_m}(\wallX_{W_m})
  -\omega^-_{W_m}(\wallX_{W_m})\bigr].
  \label{eq:first-moment-identity}
\end{equation}
Third, tightness bounds the first absolute moment: if (LR3) held, pick $M_0$
with $\sup_m\sum_{\abs{\nu}>M_0}(1+\abs{\nu})p_{W_m}(\nu)\le1$; splitting the
sum at $M_0$ and using total mass $1$ gives
$\sum_\nu\abs{\nu}p_{W_m}(\nu)\le M_0+1$ for every $m$. Combining
\eqref{eq:first-moment-identity} with (M-ESC) makes the left side diverge
linearly in $\abs{W_m}$, contradicting that uniform bound; the quantitative
form
$\sum_{\abs{\nu}>M}(1+\abs{\nu})p_{W_m}(\nu)\ge
\bigl|\sum_\nu\nu p_{W_m}(\nu)\bigr|-M$ then gives (b), and (c) is D27's
definition of $\delta x$ read against (a). Reading the third ingredient
forwards instead of backwards gives \eqref{eq:positive-transport-bound}.

Non-circularity matters here and is checked: the first-moment identity is
consumed at a step whose own leaf justifications name only (LR1), (LR2), and
the exact finite-time mean formula --- not the (LR3)-conditional statement of
the spectral memory-index theorem --- so the proof by contradiction against
(LR3) does not assume what it refutes.

Finally, the mechanism. If in the incoming Ces\`aro states the window lies
entirely in one tail in the mean, while in the outgoing states the kink has
transited beyond $W_m$ with weight $p_{\mathrm{tr}}$ and stayed with weight
$1-p_{\mathrm{tr}}$, then the definition
$\wallX_W=a-1+(1/2s)\sum_{x\in W}(S^z_x+s)$ gives
$\omega^-_{W_m}(\wallX_{W_m})=a_m-1+o(\abs{W_m})$ and
$\omega^+_{W_m}(\wallX_{W_m})=a_m-1+p_{\mathrm{tr}}\abs{W_m}+o(\abs{W_m})$, so
$\theta_{\mathrm{tr}}=p_{\mathrm{tr}}$. Only the \emph{mean} tail densities
enter: no variance hypothesis, no channel projections, no wave operators, no
sign-definite velocity classes, no incoming-concentration step.
\end{proof}

\begin{remark}[The dichotomy, and what it costs the memory ledger]
\label{rem:dichotomy}
Taken together, \cref{subsec:eps-theorem} and \cref{thm:escape-obstruction}
partition the reach of the \emph{unsubtracted} window charge. On
confined-charge sharp-tail states, where (K-TAIL) and (K-Q) hold and
$\theta_{\mathrm{tr}}=0$, the $\varepsilon$-version is available and
\cref{cor:ordered-memory-zero} makes the ordered memory exactly zero. On
mean-transiting states, where $\theta_{\mathrm{tr}}>0$, clause (LR3) fails and
no ordered outcome measure exists. Neither side of the dichotomy carries a
nonzero ordered memory for the unsubtracted charge. A two-atom ledger with
$\delta x\neq0$ therefore requires a \emph{subtracted} or co-moving
observable --- the leg subtraction of the conservation-trap paragraph, a
definition-level move which is proposed but not made here.

The two-atom shape itself is real as a mechanism, and is recorded as evidence
for the arithmetic $\theta_{\mathrm{tr}}=p_{\mathrm{tr}}$ rather than
re-claimed as a theorem: a critic reproduced it to four digits on a scattering
model (weights $0.5567$ and $0.4433$, incoming-concentration defect
$2.4\times10^{-9}$, tail $(1+513)\times0.4433=227.86$ against a measured
$227.867$), and a checker gate re-certifies the configuration on a barrier
model.
\end{remark}

\provenance{ACE-LD-obst-prime}{\texttt{theory/ace-ld.md} step (1.5) (PROPOSITION ACE-LD-obst$'$), with the mechanism remark at (1.5).(2.5) and the fences at (1.5).(2.6)}{\texttt{theory/checks/ace\_ld\_check.py} gate LD-C5 (incoming concentration $\le3.2\times10^{-6}$ at every window, and the exact two-atom support $\{\nu=0,\;\nu=-\abs{W}\}$ including the $\nu=0$ atom, with weights matching independently measured $\abs{r}^2=0.2704$, $\abs{t}^2=0.7296$); gate LD-C6 (bounded-transport contrast at protocol time $t_+=4$: mass at $\abs{\nu}>3$ below $9.1\times10^{-13}$, first moment $\le0.453$, with the bound $1.2$ held across a recorded sweep $t_+\in\{4,20,40,200\}$ whose maximum is $1.0712$; red-armed by \texttt{--red-c6-moving} at $12.4024>1.2$ and \texttt{--red-c6-weaktransit} at mass $3.742\times10^{-2}$). \textbf{Neither certificate exhibits an (M-ESC) state, and no gate bears on claims (a)--(c), which are proved, not tested}; the residual caveat is that the sweep gate itself is reached by no registered mutation and is green-side evidence only}{\texttt{ace-ld-r3.md} (M2 adopted verbatim, M3), \texttt{ace-ld-r4.md} \S0(ii),(iv) (M3 repaired, M1 residual caveat), \texttt{ace-ld-r2.md}; no-witness confirmed independently in \texttt{lr-d16-r1.md} M2(b)1 and \texttt{lr-d16-r2.md} M4(b)3}

\subsection{Composition: mean escape excludes bounded escaped content}
\label{subsec:m-esc-nr}

The two conditional theorems above point in opposite directions, and composing
them is pure propositional logic.

\begin{proposition}[Mean escape and bounded escaped content are incompatible]
\label{thm:m-esc-nr}
\statusSketch\;
Under the joint hypotheses of \cref{thm:escape-obstruction} and
\cref{thm:nr-transfer} --- the same vector $\Psi$, cut, D27 sequence,
fixed-window laws, and an arbitrary padded exhaustion --- (M-ESC) implies
$\neg$(NR): at least one of the two bounds \eqref{eq:nr-clause-1},
\eqref{eq:nr-clause-2} is infinite. This holds exhaustion by exhaustion.
\end{proposition}

\begin{scope}
\label{scope:m-esc-nr}
The status is a sketch that is a candidate for promotion: the argument is a
pure composition of two exact conditional statements, and review and
adjudication are pending. Not claimed: \emph{which} clause of (NR) fails; any
model or state realising (M-ESC); a counterexample in the exactly solvable
model; or necessity of (NR) outside this named route.
\end{scope}

\begin{proof}[Idea of proof]
\Cref{thm:escape-obstruction} gives (M-ESC) $\Rightarrow$ $\neg$(LR3), while
\cref{thm:nr-transfer} gives (NR) $\Rightarrow$ (LR3). Assuming both (M-ESC)
and (NR) yields (LR3) and $\neg$(LR3). No new computation is involved, and
neither source checker exhibits an (M-ESC) witness.
\end{proof}

\provenance{M-ESC-NR}{\texttt{theory/lr-d16.md} step (1.5).(2.9), THEOREM LRD-MESC-NR; composed from \texttt{theory/ace-ld.md} (1.5) and \texttt{theory/lr-d16.md} (1.5).(2.1)--(2.7)}{No new computation: propositional composition of two exact conditional rows; neither source checker exhibits an (M-ESC) witness}{\texttt{blitz-m-esc-nr-r1.md} (0F/0M/0m)}

\subsection{Kink--magnon wave operators from exact ansatz band data}
\label{subsec:wave-operators}

The escape analysis above presumes a channel decomposition. This subsection
records the one construction of such a decomposition that the campaign has
proved --- conditionally, on an explicitly displayed and unverified
localization hypothesis. The general two-magnon theory is in
\cref{sec:two-magnon} and the abstract wave-operator layer in
\cref{sec:wave-operators}; here we state the kink--magnon case and its
hypothesis package in full.

\begin{definition}[Exact fixed-packet kink--magnon band data]
\label{def:d28}
Assume H-MQG(1)--(3). In addition:
\begin{enumerate}
\item \emph{(Covariant kink realization.)} The kink folium has a covariant
positive-energy Hilbert realization
$(\Hilb_{\alpha\beta},\pi_{\alpha\beta},U_{\alpha\beta})$ implementing the
commuting infinite-volume time and lattice translations. This is extra
structure beyond the state-set definition. The vacuum GNS triples are
$(\Hilb_\gamma,\pi_\gamma,\Omega_\gamma)$ for $\gamma\in\{\alpha,\beta\}$, with
$H\Omega_\gamma=0$; there is \emph{no} vacuum vector in
$\Hilb_{\alpha\beta}$.
\item \emph{(Exact kink band.)} Its Hamiltonian has an isolated
finite-multiplicity kink band with dispersion $E_K\in C^2(\mathbb{T})$, matrix
valued on the finite frame index. A fixed Gram-normalised,
translation-covariant finite-core ansatz frame gives an \textbf{exact} band map
$\Gamma_K$ satisfying $H_{\alpha\beta}\Gamma_K=\Gamma_KE_K$; it is not merely
a Rayleigh--Ritz approximation.
\item \emph{(Exact magnon bands with momentum-filtered creators.)} Each
relevant vacuum tail $\gamma$ has an isolated selected magnon band
$\omega_\gamma$ (written $\omega$ when symmetry relates the tails), gap
$\Delta_M:=\min_k\omega(k)>0$, and Gram-normalised \textbf{exact} band map
$\Gamma_{M,\gamma}$. A spacetime-Schwartz energy--momentum filter of a local
observable gives almost-local creators $a_{\gamma,b}(n)$ with strictly positive
energy transfer and with the \emph{momentum-filtered normalisation}
\begin{equation}
  a_{\gamma,b}(n)\Omega_\gamma
  \;=\;\Gamma_{M,\gamma}\bigl(\chi_\gamma e_n\otimes e_b\bigr),
  \qquad e_n(k):=e^{-ikn},
  \label{eq:filtered-creator}
\end{equation}
for a fixed filter $\chi_\gamma\in C_c^\infty(\mathbb{T})$; translation
covariance makes $\chi_\gamma$ independent of $n$. This \emph{equality}, not
merely nonzero overlap with the band, is part of the exactness hypothesis. A
position-delta normalisation, the case $\chi_\gamma\equiv1$, is not assumed and
is in fact incompatible with a compactly supported energy--momentum filter.
\item \emph{(Fixed packets, separated velocities, filter compatibility.)}
Packet amplitudes are finite sums of $C_c^\infty$ products and the resulting
physical packet states lie in the $\ell^1$ class D17. The dispersions are
$C^\infty$ on open neighbourhoods $U_K,U_\gamma$ of their packet momentum
supports (global $C^2$ is retained), their kink and magnon velocity supports
have distance $\varepsilon_v>0$, and their signs are: incoming-left
$v_M-v_K\ge\varepsilon_v$ at $t\to-\infty$; outgoing-left
$v_M-v_K\le-\varepsilon_v$; outgoing-right $v_M-v_K\ge\varepsilon_v$ at
$t\to+\infty$. This is a fixed-packet condition and excludes equal velocities
and the soft endpoint. The filters of clause 3 and the kink filter
$\chi_K\in C_c^\infty(\mathbb{T})$ satisfy $\chi_\bullet\equiv1$ on the
corresponding packet momentum support with
$\supp\chi_\bullet\subset U_\bullet$, and companions
$\widetilde{\chi}_\bullet\in C_c^\infty(\mathbb{T})$ satisfy
$\widetilde{\chi}_\bullet\equiv1$ on $\supp\chi_\bullet$ with
$\supp\widetilde{\chi}_\bullet\subset U_\bullet$. The channel Hamiltonian
symbols $h_K:=E_K\widetilde{\chi}_K$ and
$h_\gamma:=\omega_\gamma\widetilde{\chi}_\gamma$ are then
$C^\infty(\mathbb{T})$ and agree with the dispersions on the packet supports.
\item \emph{(Two-cluster factorization through the vacuum tails.)} Write
$\kappa_a^{(0)}(x):=\Gamma_K(e_x\otimes e_a)$ for the unfiltered exact kink
Wannier vectors. There are $C_{\mathrm{cl}}<\infty$ and
$\lamt\in\bigl(\max(\lambda_{E_\alpha},\lambda_{E_\beta}),1\bigr)$ such that
for every $\gamma\in\{\alpha,\beta\}$, all frame indices $a,a'$, all
$x,x'\in\Z$, every $r\ge1$, and every pair $A\in\alg_{\Lambda_\alpha}$,
$B\in\alg_{\Lambda_\beta}$ with
$\Lambda_\alpha\subset(-\infty,\min(x,x')-r]$ on the $\alpha$ side and
$\Lambda_\beta\subset[\max(x,x')+r,\infty)$ on the $\beta$ side (either factor
may be $\mathbb{1}$),
\begin{equation}
  \Bigl|\bigl\langle\kappa_a^{(0)}(x),\,AB\,\kappa_{a'}^{(0)}(x')\bigr\rangle
  -\omega_\alpha(A)\,\omega_\beta(B)\,
   \bigl\langle\kappa_a^{(0)}(x),\kappa_{a'}^{(0)}(x')\bigr\rangle\Bigr|
  \;\le\;C_{\mathrm{cl}}\norm{A}\,\norm{B}\,\lamt^{\,r}.
  \label{eq:d28-c}
\end{equation}
The one-tail case is $B=\mathbb{1}$ (or $A=\mathbb{1}$); the genuinely
two-sided case is used at exactly one step. The constants depend on the model
and the frame only --- \emph{not} on $A,B,\Lambda_\alpha,\Lambda_\beta,r,x,x',
a,a'$, and in particular not on any time parameter, the positions being
packet-driven and time dependent downstream. This is the localization input
that turns exact band equations into a two-cluster estimate; isolated
eigenvalues alone do not imply it. It is the sole replacement for the vacuum
clustering apparatus of the source literature, and no four-point analogue is
assumed.
\end{enumerate}
The asymptotic channel Hamiltonian is multiplication by
$E_K(p)+\omega_\gamma(k)$. Smooth packet domains with the three velocity signs
are denoted $D_-^L,D_+^L,D_+^T$, and $P_L,P_T$ are the two outgoing direct-sum
projections. We refer to \eqref{eq:d28-c} as the \emph{two-cluster
inequality}.
\end{definition}

\begin{remark}[A regime fence, which is not a hypothesis]
\label{rem:d28-r}
Put $K_-:=\min E_K$, $K_+:=\max E_K$, $W_K:=K_+-K_-$, and let $I_{\mathrm{in}}$
be the range of $E_K(p)+\omega(k)$ on the incoming packet support. Same-band
kink absorption is excluded by
$d_{\mathrm{abs}}:=\dist(I_{\mathrm{in}},[K_-,K_+])>0$, for which the
packet-independent sufficient condition is $\Delta_M>W_K$. With
$\Delta_r:=\min_k\omega_r(k)$ for every other known one-particle vacuum band
and $K_*:=\inf_{j\ge2,p}E_{K,j}(p)$ for every other known kink band, one asks
$\sup I_{\mathrm{in}}+\eta_{\mathrm{inel}}<\Theta_{\mathrm{inel}}$ where
$\Theta_{\mathrm{inel}}:=\min\{K_-+2\Delta_M,\;K_*,\;\inf_r(K_-+\Delta_r)\}$
and $\eta_{\mathrm{inel}}>0$, absent entries being $+\infty$; the three entries
respectively exclude an additional magnon, a higher-kink absorption channel,
and a different one-particle channel. An uncomputed wall--magnon bound band is
not excluded. \textbf{These inequalities are used nowhere in the proof.} Their
role is to make the \emph{exactness} hypotheses plausible --- an isolated exact
band is not to be expected when a competing channel is open --- so they belong
to the justification of \cref{def:d28}, not to its inferential content. Adding
or deleting them changes no step.
\end{remark}

\begin{theorem}[Fixed-packet kink--magnon wave operators exist and are isometric]
\label{thm:ac-ex}
\statusProvedConditional\;
Proved only as the following conditional implication. Assume H-MQG(1)--(3) and
the whole of \cref{def:d28} --- exact kink and magnon band maps, covariant
kink-sector realization, momentum-filtered creator normalisation, velocity
separation $\varepsilon_v>0$, and the displayed \emph{two-sided} two-cluster
inequality \eqref{eq:d28-c}. Then, on the fixed packet domains:
\begin{enumerate}
\item[(ACE.1)] the Cook limits $W_-^L$, $W_+^L$ and $W_+^T$ exist;
\item[(ACE.2)] $W_-^L$ and $W_+:=W_+^L\oplus W_+^T$ are isometries, and
$N_T^{\mathrm{ex}}:=W_+P_TW_+^*$ is an orthogonal projection on the
constructed out-space $\Hilb_{\mathrm{out}}^{\mathrm{ex}}:=\ran W_+$ --- the
transmitted-channel number of D18 restricted to that range, with nothing
asserted on its orthogonal complement.
\end{enumerate}
A third conclusion is \textbf{triply conditional}, on \emph{all three} of: the
charge assignment $q_{\mathrm{in}}=q_L=-1$, $q_T=+1$; the named local-decay
hypothesis (AD3-ex) of \cref{conj:ad3-ex}, which this construction does not
prove; and the existence of the ordered-limit outcome measure for the vector in
question, which \cref{def:d28} does not supply. Under all three:
\begin{enumerate}
\item[(ACE.3)] for a normalized event vector in
$\ran W_-^L\cap\Hilb_{\mathrm{out}}^{\mathrm{ex}}$ \emph{for which that outcome
measure exists}, the constructed-channel part of the measure has support
\emph{contained in} $\{0,2\}$ (either weight may vanish),
\begin{equation}
  p_\nu^{\mathrm{ex}}
  =\norm{P_LW_+^*\Psi}^2\delta_{\nu,0}
  +\norm{P_TW_+^*\Psi}^2\delta_{\nu,2},
  \qquad
  p_2^{\mathrm{ex}}=\braket{\Psi,N_T^{\mathrm{ex}}\Psi},
  \label{eq:ace3-ledger}
\end{equation}
and the ledger gives $\delta x_{\mathrm{ex}}=-p_2^{\mathrm{ex}}/s$.
\end{enumerate}
The order of limits is frozen: infinite-volume dynamics first, then the
wave-operator time limit, then any fixed-window limit, and the spatial
exhaustion last. No statement is made at zero momentum or after a packet-soft
limit.
\end{theorem}

\begin{scope}
\label{scope:ac-ex}
\textbf{The two-cluster inequality \eqref{eq:d28-c} is the load-bearing
hypothesis and is unverified on any model.} The threshold inequalities of
\cref{rem:d28-r} are used in no step. No completeness, no implication from raw
band data, no bound-state exclusion, no threshold-inequality use, and no soft
limit is claimed. Nothing is asserted about the orthogonal complement of the
constructed out-space.
\end{scope}

\begin{proof}[Idea of proof]
The construction is a Cook argument with the vacuum clustering apparatus
replaced by the single displayed cut \eqref{eq:d28-c}. One fixes a smooth
monotone profile $\theta$ with $\theta'$ supported in $[-w,w]$ and
$\norm{\theta'}_\infty\le C/w$, sets $\theta_T:=\theta$ and
$\theta_L:=1-\theta$, and defines the kink--magnon precursor identification
\begin{equation}
  I_cF\;:=\;\sum_{x,y,a,b}\theta_c(y-x)\,F^{ab}(x,y)\,
  a_{c,b}(y)\,\kappa_a(x),
  \label{eq:precursor-identification}
\end{equation}
with the filtered Wannier frame $\kappa_a(x)=\Gamma_K(\chi_Ke_x\otimes e_a)$.
Both profiles are fixed, independent of $t$ and of the packet. The free channel
Hamiltonian is multiplication by $E_K(p)+\omega_\gamma(k)$, which on the packet
domain coincides with convolution by the filtered kernels, because the packet
momentum supports lie where the filters are identically one. Commuting the
creator past the Hamiltonian produces a defect that splits into an exact
band-equation term and two kernel-tail terms; velocity separation makes the
supports of the kink and magnon packets separate ballistically, so the defect
is integrable in $t$ and the Cook limits exist.

Isometry is where the two-cluster inequality is consumed. One needs a
four-cluster factorization,
\begin{equation}
  \bigl|\braket{a_{\gamma,b}(y)\kappa_a(x),\,
  a_{\gamma,b'}(y')\kappa_{a'}(x')}
  -\braket{\kappa_a(x),\kappa_{a'}(x')}\,G_\gamma^{bb'}(y,y')\bigr|
  \;\le\;C_N\langle r\rangle^{-N},
  \label{eq:four-cluster}
\end{equation}
with $G_\gamma$ the exact magnon Gram form. This is \emph{derived}, not
assumed: the product $a_{\gamma,b}(y)^\dagger a_{\gamma,b'}(y')$ is almost
local about $\{y,y'\}$, so it may be truncated to a one-sided region at
distance $\gtrsim r$ with rapidly decaying error, whereupon the one-tail case
of \eqref{eq:d28-c} factors it and the vacuum evaluation is exactly the magnon
Gram form by the filtered-creator normalisation
\eqref{eq:filtered-creator}. The genuinely two-sided case of \eqref{eq:d28-c}
is used at one step only. The isometry and the projection property of
$N_T^{\mathrm{ex}}$ follow on the constructed range, with no completeness
statement anywhere.
\end{proof}

\provenance{AC-EX; D28}{\texttt{theory/ansatz-scattering.md}: the theorem at \S0, the almost-locality and defect estimates at step (1.3), velocity separation at (1.4), Cook integrability at (1.5), isometry and channel number at (1.6), and the conditional ledger at (1.7)}{\texttt{theory/checks/ansatz\_scattering\_check.py}, gates ACE-C1a/C1b/C2a/C2b (green exit $0$; \texttt{--red}, \texttt{--red-slow-kernel}, \texttt{--red-equal-velocity}, \texttt{--red-absorption} each exit $1$; all five re-run independently). Kernel decay is red-certified load-bearing on the majorant of step (1.5).(2.4)(i), not on the literal $\ell^2$ defect, which is insensitive to it and whose mutation is registered dead; \texttt{--red-absorption} is a fence regression guard and is \emph{not} a red test of this theorem. \textbf{Neither certificate touches the two-cluster inequality \eqref{eq:d28-c}}}{\texttt{memory-index-r3.md} \S4 (amended sentence; the earlier hold on this row is lifted) and \S0 V6, \S1}

\subsection{The open local-decay hypothesis}
\label{subsec:ad3-ex}

What separates the conditional ledger (ACE.3) from a theorem is one named
hypothesis, and it is worth displaying exactly as it stands.

\begin{conjecture}[Local decay on the constructed channels]
\label{conj:ad3-ex}
\statusConjecture\;
Let $\Psi=W_+(F_L,F_T)$ and $\Psi_t:=e^{-itH}\Psi$. For every fixed window $W$
\textbf{containing the kink core region}, with $E_W(\cdot)$ the spectral
measure of the regularised window charge and $P_{\mathrm{ch}}$,
$\mathrm{ch}\in\{L,T\}$, the constructed channel projections,
\begin{equation}
  \lim_{t\to+\infty}
  \norm{E_W(\{q_{\mathrm{ch}}\})\Psi_t-P_{\mathrm{ch}}\Psi_t}=0,
  \label{eq:ad3-ex}
\end{equation}
the limit being taken \emph{before} $W\uparrow\Z$.
\end{conjecture}

\begin{scope}
\label{scope:ad3-ex}
The order of limits is part of the statement: time first at fixed $W$, then the
window exhaustion. The hypothesis is the corresponding clause of the general
asymptotic-decay package D18, restricted to the constructed channels;
\cref{def:d28} deliberately does not contain it. No part of it is proved. What
\emph{is} proved from it is the vanishing of the inter-channel cross terms, at
step (1.7).(2.2) of the construction. The gloss ``the outgoing
boundary-straddling charge on $\partial W$ vanishes in the ordered limit''
describes the \emph{missing mechanism}, not the hypothesis: the free leg and
its dressing do leave any fixed window at rate $O(\abs{t}^{-N})$, but
converting that into convergence of the window-charge \emph{spectral
projection} requires control of the boundary-straddling charge that no step of
the construction supplies. \Cref{thm:spectral-diagonality} is the abstract
statement of what such control would buy, and \cref{thm:eps-theorem} is the
$\varepsilon$-weakened version available on the confined-tail class --- which,
by \cref{cor:no-escaping-leg}, is disjoint from the class this hypothesis is
about.
\end{scope}

\provenance{AD3-ex}{Displayed in \texttt{theory/ansatz-scattering.md} step (1.7), in the ASSUME block; consumed at (1.7).(2.2)}{---}{---}
